\begin{musica}{Engenharia}{Paulino Gomes Júnior}{Tuna de Engenharia da Universidade do Porto}
\nota{Introdução: D E Am}

Engenha\ac{Am}{ria},
Curso que me envai\ac{E}{de}ce
É ale\ac{E7}{gria},
Chama imensa que engran\ac{Am}{de}ce. \ac{A7}{}__
Faz de nós conquista\ac{Dm}{do}res, \ac{B7}{}__
Génios, bravos luta\ac{E}{do}res.
\ac{Dm}{Faz} do Porto a capi\ac{Am}{tal}
Do estu\ac{B7}{dan}te em \ac{E}{Por}tu\ac{A}{gal}.

\nota{Refrão:}
Engenha\ac{A}{ria},
É grande \ac{F#m}{luz} d’acade\ac{Bm}{mia}
É facul\ac{E}{da}de onde há sa\ac{E7}{ber} e ami\ac{A}{za}de,
Que o \ac{Dm}{Sol} lá do céu
Risonho vem bei\ac{A}{jar}
Com um orgulho muito \ac{Bm}{seu}, \ac{E}{}_
E as capas ondu\ac{A}{lan}tes
\ac{F#m}{Põe} no peito a vi\ac{Bm}{brar}
\ac{E}{O} coração dos estu\ac{A}{dan}tes.

\nota{Instrumental: A F#m Bm E E7 A}

Que o Sol lá do céu Risonho vem beijar
Com um orgulho muito seu, E as capas ondulantes
Põe no peito a vibrar
O coração dos estudantes.

Põe no peito a vi\ac{Dm}{brar}
O coração dos estu\ac{A}{dan}tes.

\posicoes{D: xx0232, Dm: x00231, E: 022100, E7: 022130, F#m: 244222, A: x02220,
  Am: x02210, A7: x02020, Bm: x24432, B7: x24242}
\end{musica}
